\chapter{Global and Interior Gradient Estimates}

Let
\begin{equation}\tag{15.1}
Qu=a^{ij}(x,u,Du)D_{ij}u+b(x,u,Du)=0
\end{equation}
be a quasilinear elliptic equation on a domain $\Omega\subset\mathbb R^n$.
The estimates below control $|Du|$ globally from boundary data and locally from
the oscillation and distance to the boundary.

\section*{15.1. A Maximum Principle for the Gradient}

Assume first that the principal coefficients admit
\begin{equation}\tag{15.2}
a^{ij}(x,z,p)=a_*^{ij}(p)
+
\frac12\bigl(p_ic_j(x,z,p)+c_i(x,z,p)p_j\bigr),
\end{equation}
where $a_*^{ij}\in C^1(\mathbb R^n)$,
$c_i\in C^1(\Omega\times\mathbb R\times\mathbb R^n)$, and
$[a_*^{ij}]$ is nonnegative and symmetric.


(i) If \(a^{ij}=a^{ij}(p)\), take \(a_*^{ij}=a^{ij}\) and \(c_i=0\).

(ii) The Euler--Lagrange equation corresponding to a multiple integral of the form

\begin{equation}
\int_{\Omega} F(x,u,|Du|)\,dx,
\tag{15.3}
\end{equation}

where \(F\in C^2\), \(D_tF\ne0\), and \(t=|p|\), has equation

\begin{equation}
\Delta u+\{(|Du|D_{tt}F/D_tF)-1\}|Du|^{-2}D_i u D_j u D_{ij}u
+ (|Du|^2D_{tz}F + D_i u D_{tx_i}F - |Du|D_zF)/D_tF=0,
\tag{15.4}
\end{equation}

so that the decomposition (15.2) is valid with

\[
a_*^{ij}=\delta^{ij},\qquad c_i=\{(|p|D_{tt}F/D_tF)-1\}\,|p|^{-2}p_i
\]

provided \(c_i\in C^1\).

(iii) In the special case of two variables we can write

\begin{equation}
 a^{ij}=(\Tr-\Estar)\delta^{ij}+\tfrac12(p_i d_j+d_i p_j)
\tag{15.5}
\end{equation}

where

\[
\Tr=\operatorname{trace}[a^{ij}]=a^{11}+a^{22}
\]
\[
\Estar=\mathcal{E}/|p|^2=a^{ij}p_ip_j/|p|^2,
\]
\[
d_1=[(a^{11}-a^{22})p_1+2a^{12}p_2]/|p|^2,
\]
\[
d_2=[2a^{12}p_1+(a^{22}-a^{11})p_2]/|p|^2.
\]

For \(n=2\), (15.1) is equivalent to

\begin{equation}
\Delta u+c_iD_j uD_{ij}u+b^*=0
\tag{15.6}
\end{equation}

where \(c_i=d_i/(\Tr-\Estar)\), \(b^*=b/(\Tr-\Estar)\), and
\(a_*^{ij}=\delta^{ij}\).

Set \(v=|Du|^2\). By (15.2), equation (15.1) becomes

\[
a_*^{ij}(Du)D_{ij}u+\tfrac12 c_i(x,u,Du)D_iv+b(x,u,Du)=0.
\]




For \(u\in C^3(\Omega)\), differentiate, multiply by \(D_ku\), and sum over \(k\):

\begin{equation}
\tag{15.7}
-2a_*^{ij}D_{ik}uD_{jk}u+a_*^{ij}D_{ij}v+\bigl(D_{p_i}a_*^{ij}D_{ij}u+D_{p_i}b+D_kuD_kc_i\bigr)D_iv+2\delta bv=0
\end{equation}

Here

\begin{equation}
\tag{15.8}
\delta=D_z+|p|^{-2}p_iD_{x_i}.
\end{equation}

Schwarz's inequality gives

\begin{equation}
\tag{15.9}
a_*^{ij}D_{ik}uD_{jk}u\ge (a_*^{ij}D_{ij}u)^2/\operatorname{trace}[a_*^{ij}]
=(b+\tfrac12c_iD_iv)^2/\operatorname{trace}[a_*^{ij}].
\end{equation}

Thus

\begin{equation}
\tag{15.10}
a^{ij}D_{ij}v+B_iD_iv\ge 2\left(\frac{b^2}{\mathcal{J}_*}-v\,\delta b\right)
\end{equation}

where $\mathcal{J}_*=\operatorname{trace}[a_*^{ij}]$ and

\[
B_i=D_{p_i}a_*^{ij}D_{ij}u+D_{p_i}b+D_kuD_kc_i-2\mathcal{J}_*^{-1}bc_i-\frac12\mathcal{J}_*^{-1}c_ic_jD_jv.
\]

Assume

\begin{equation}
\tag{15.11}
|p|^2\delta b\le \frac{b^2}{\mathcal{J}_*}
\end{equation}

on \(\Omega\times\mathbb R\times\mathbb R^n\). For \(C^2\) solutions, the
divergence form of (15.7) has weak identity
\begin{equation}
\tag{15.12}
-\int_\Omega \bigl\{(D_{p_i}a_*^{ij}D_{ij}u-D_ja^{ij}+D_kuD_kc_i)D_iv\bigr\}\eta\,dx
+\int_\Omega 2b\eta\,\Delta u\,dx=0
\end{equation}



for \(\eta\in C_0^1(\Omega)\). Integration by parts yields

\begin{equation}
\tag{15.13}
D_i(a^{ij}D_j v)+(B_i-D_ja^{ij})D_iv \ge 2(b^2/\mathcal{J}_* - v\,\delta b)
\end{equation}

for \(v\in C^1(\bar\Omega)\).

\medskip

\begin{theorem}[Gradient maximum principle]\label{gilbarg-trudinger.thm.15.1.gradient-maximum-principle}\dcref{ch15:15.1}\group{ch15-global-interior-gradient}
Let $u \in C^2(\Omega) \cap C^1(\bar{\Omega})$ satisfy equation (15.1) in the bounded
domain $\Omega$ and suppose that $Q$ is elliptic in $\Omega$ with coefficients satisfying (15.2) and
(15.11). Then

\begin{equation}
\tag{15.14}
\sup_{\Omega} |Du| = \sup_{\partial\Omega} |Du|
\end{equation}
\end{theorem}
\begin{proof}
\sketch For $v=|Du|^2$, differentiation of (15.1), (15.2), and Schwarz's inequality give
(15.10); for $C^2$ solutions the approximation argument yields the weak form (15.13).
Condition (15.11) makes the right-hand side nonnegative, so the maximum principle gives (15.14).
\end{proof}

The hypotheses need only hold for \(|z|\leq\sup_\Omega|u|\). If they hold for
\(|p|\ge L\), then

\begin{equation}
\tag{15.15}
\sup_{\Omega} |Du| \leq \max \left\{ \sup_{\partial\Omega} |Du|,\, L \right\}.
\end{equation}

This follows by applying Theorem 15.1 on \(\Omega_L=\{x\in\Omega:|Du|>L\}\).
For \(a_*^{ij}=\delta^{ij}\), (15.11) becomes

\begin{equation}
\tag{15.16}
|p|^2\,\delta b \leq \frac{b^2}{n}
\end{equation}

In particular, \(b=0\) gives a gradient maximum principle.

\bigskip

\section*{15.2. The General Case}

\medskip

Set \(m=\inf_\Omega u\), \(M=\sup_\Omega u\), and choose an increasing
\(\psi\in C^3[\bar m,\bar M]\) with \(m=\psi(\bar m)\), \(M=\psi(\bar M)\).
Write \(u=\psi(\bar u)\), so

\[
D_i u=\psi' D_i \bar u
\]
\[
D_{ij}u=\psi'' D_i \bar u D_j \bar u+\psi' D_{ij}\bar u,
\]

and

\begin{equation}
\psi' a^{ij}(x,u,Du)D_{ij}\bar u+b(x,u,Du)+\frac{\psi''}{(\psi')^2}\mathcal E(x,u,Du)=0.
\tag{15.17}
\end{equation}

With \(v=|Du|^2\), \(\bar v=|D\bar u|^2\), differentiation of (15.17) gives

\begin{align}
&-a^{ij}D_{ik}\bar u D_{jk}\bar u+\frac12 a^{ij}D_{ij}\bar v+\frac{1}{2(\psi')^2}\Bigl(\psi' D_{p_l}a^{ij}D_{ij}\bar u+D_{p_l}b
\notag\\
&\qquad+\frac{\psi''}{(\psi')^2}D_{p_l}\mathcal E\Bigr)D_l v+\Bigl(\psi' \delta a^{ij}D_{ij}\bar u+\delta b+\frac{\psi''}{(\psi')^2}\delta \mathcal E\Bigr)\bar v
\tag{15.18}\\
&\qquad-\frac{\psi''}{(\psi')^2}\Bigl(b+\frac{\psi''}{(\psi')^2}\mathcal E\Bigr)\bar v-\frac{1}{\psi'}\Bigl[\frac{\psi''}{(\psi')^2}\Bigr]'\delta \bar v=0.
\notag
\end{align}

Define

\begin{equation}
\delta=p_iD_{p_i}
\tag{15.19}
\end{equation}

and

\begin{equation}
\omega=\frac{\psi''}{(\psi')^2}\in C^1[\bar m,\bar M],
\tag{15.20}
\end{equation}

Using

\[
D_l v=2\psi''\bar v D_l\bar u+(\psi')^2D_l\bar v,
\]

we obtain

\begin{align}
&-a^{ij}D_{ik}\bar u D_{jk}\bar u+\frac12 a^{ij}D_{ij}\bar v+\frac12\Bigl(\psi' D_{p_i}a^{ik}D_{jk}\bar u+D_{p_i}b+\omega D_{p_i}\mathcal E\Bigr)D_i \bar v
\notag\\
&\qquad+\psi'\bar v(\omega \delta a^{ij}+\delta a^{ij})D_{ij}\bar u
\tag{15.21}\\
&\qquad+\Bigl\{\frac{\omega'}{\psi'}\mathcal E+\omega^2(\delta-1)\mathcal E+\omega(\delta \mathcal E+(\delta-1)b)+\delta b\Bigr\}\bar v=0.
\notag
\end{align}

For scalar multipliers \(r,s\), combine (15.17) and (15.21):

\begin{equation}
\begin{aligned}
&-a^{ij}D_{ik}\bar{u}D_{jk}\bar{u}+\tfrac12a^{ij}D_{ij}\bar{v}
+\tfrac12\bigl(\psi' D_{p_i}a^{jk}D_{jk}\bar{u}+D_{p_i}b+\omega D_{p_i}\mathcal{E}\bigr)D_i\bar{v}\\
&\qquad+\psi'\bar{v}\bigl(\omega(\delta+r+1)a^{ij}+(\delta+s)a^{ij}\bigr)D_{ij}\bar{u}\\
&\qquad+\left\{\frac{\omega'}{\psi'}\mathcal{E}+\omega^2(\delta+r)\mathcal{E}
+\omega\bigl((\delta+s)\mathcal{E}+(\delta+r)b\bigr)+(\delta+s)b\right\}\bar{v}=0.
\end{aligned}
\tag{15.22}
\end{equation}

At this point it is convenient to write the principal coefficients $a^{ij}$ in the form

\begin{equation}
a^{ij}(x,z,p)=a^{ij}_{*}(x,z,p)+\tfrac12\bigl[p_ic_{j}(x,z,p)+c_i(x,z,p)p_j\bigr]
\tag{15.23}
\end{equation}

where \(a_*^{ij},c_i\in C^1\) and \([a_*^{ij}]\) is positive and symmetric.
Substitution in (15.22) gives

\begin{equation}
\begin{aligned}
&-a^{ij}_{*}D_{ik}\bar{u}D_{jk}\bar{u}+\tfrac12a^{ij}D_{ij}\bar{v}
+\tfrac12\bigl(\psi' D_{p_i}a^{jk}D_{jk}\bar{u}-\psi'c_jD_{ij}\bar{u}\\
&\qquad+\bar{v}\bigl[\omega(\delta+r+1)+\delta+s+1\bigr]c_i+D_{p_i}b+\omega D_{p_i}\mathcal{E}\bigr)D_i\bar{v}\\
&\qquad+\psi'\bar{v}\bigl[\omega(\delta+r+1)a^{ij}+(\delta+s)a^{ij}_{*}\bigr]D_{ij}\bar{u}\\
&\qquad+\left\{\frac{\omega'}{\psi'}\mathcal{E}+\omega^2(\delta+r)\mathcal{E}
+\omega\bigl[(\delta+s)\mathcal{E}+(\delta+r)b\bigr]+(\delta+s)b\right\}\bar{v}=0.
\end{aligned}
\tag{15.24}
\end{equation}

By Cauchy's inequality (7.6),

\[
\psi'\bar{v}\bigl[\omega(\delta+r+1)a^{ij}+(\delta+s)a^{ij}_{*}\bigr]D_{ij}\bar{u}
\]
\[
\le \lambda_{*}\sum\left|D_{ij}\bar{u}\right|^2+\frac{\bar{v}}{4\lambda_{*}}
\sum\left[\omega(\delta+r+1)+\delta+s\right]\left[a^{ij}_{*}\right]^2
\]
\[
\le a^{ij}_{*}D_{ik}\bar{u}D_{jk}\bar{u}+\frac{\bar{v}}{4\lambda_{*}}
\sum\left[\omega(\delta+r+1)+\delta+s\right]\left[a^{ij}_{*}\right]^2
\]

where \(\lambda_*\) is the minimum eigenvalue of \([a_*^{ij}]\). Hence

\begin{equation}
a^{ij}D_{ij}\bar{v}+B_iD_i\bar{v}+2G\mathcal{E}\bar{v}\ge 0
\tag{15.25}
\end{equation}

where

\begin{equation}
B_i=\psi'\bigl(D_{p_i}a^{jk}D_{jk}\bar{u}-c_jD_{ij}\bar{u}\bigr)+\bar{v}\bigl[\omega(\delta+r+1)+\delta+s+1\bigr]c_i
+D_{p_i}b+\omega D_{p_i}\mathcal{E},
\tag{15.26}
\end{equation}

\begin{equation}\tag{15.27}
\begin{aligned}
G&=\frac{\omega'}{\psi'}+\alpha\omega^2+\beta\omega+\gamma,\\
\alpha &= \frac{1}{\mathcal{E}_*}\left((\delta+r)\mathcal{E} + \frac{|p|^2}{4\lambda_*}\sum\bigl((\delta+r+1)a_*^{ij}\bigr)^2\right),\\
\beta &= \frac{1}{\mathcal{E}_*}\left((\delta+s)\mathcal{E}+(\delta+r)b+\frac{|p|^2}{2\lambda_*}\bigl[(\delta+r+1)a_*^{ij}\bigr]\bigl[(\delta+s)a_*^{ij}\bigr]\right),\\
\gamma &= \frac{1}{\mathcal{E}_*}\left[\frac{|p|^2}{4\lambda_*}\sum\bigl((\delta+s)a_*^{ij}\bigr)^2+(\delta+s)b\right].
\end{aligned}
\end{equation}

Choose \(\psi,r,s\) so that

\begin{equation}\tag{15.28}
G\le 0 \qquad \text{for } x\in \Omega,\qquad \text{and}\qquad |Du(x)|\ge L
\end{equation}

for sufficiently large \(L\). Then the weak maximum principle gives

\[
\sup_{\Omega_L}\bar v=\sup_{\partial\Omega_L}\bar v,\qquad
(\Omega_L=\{x\in \Omega\mid |Du(x)|\ge L\}),
\]

and hence

\begin{equation}\tag{15.29}
\sup_{\Omega}|Du|\le \max\left\{\left[\frac{\max \psi'}{\min \psi'}\right]\sup_{\partial\Omega}|Du|,\,L\right\},
\end{equation}

For \(u\in C^2(\Omega)\), the weak form is

\begin{equation}\tag{15.30}
\begin{aligned}
\int_{\Omega}\eta a^{ij}D_{ik}\bar u\,D_{jk}\bar u\,dx
&+\frac{1}{2}\int_{\Omega} a^{ij}D_i\bar v\,D_j\eta\,dx\\
+\frac{1}{2}\int_{\Omega}\left\{D_j a^{ij}D_i\bar v-\frac{1}{(\psi')^2}\bigl(\psi' D_p a^{jk}D_{jk}\bar u + D_p b + \omega D_p \mathcal{E}^{\circ}\bigr)D_i\bar v\right\}\eta\,dx\\
-\int_{\Omega}\left\{\psi'\,\delta a^{ij}D_{ij}\bar u+\delta b-\omega(b+\omega\mathcal{E}^{\circ})+\frac{\omega'}{\psi'}\,\mathcal{E}^{\circ}\right\}\bar v\eta\,dx=0
\end{aligned}
\end{equation}

for all \(\eta\ge0\), \(\eta\in C_0^1(\Omega)\). It implies the weak form

\begin{equation}\tag{15.31}
\int_{\Omega}\left\{a^{ij}D_i\bar v\,D_j\eta+(D_j a^{ij}-B_i)D_i\bar v\,\eta-2G\mathcal{E}^{\circ}\bar v\eta\right\}\,dx\le 0
\end{equation}



for all nonnegative \(\eta\in C_0^1(\Omega)\); the weak maximum principle gives
(15.29). Impose

\begin{equation}
\tag{15.32}
(\delta+r+1)a^{ij}_{*},\,(\delta+s)a^{ij}_{*}=O\!\left[\frac{\sqrt{\lambda_{*}\mathcal{E}}}{|p|}\right]
\quad\text{as }|p|\to\infty
\end{equation}

\[
\delta\mathcal{E},\,\delta\mathcal{E},\,(\delta+r)b,\,(\delta+s)b \le O(\mathcal{E}) \quad \text{as }|p|\to\infty.
\]

uniformly for \((x,z)\in\Omega\times[m,M]\), and define

\begin{equation}
\tag{15.33}
a,\,b,\,c = \overline{\lim}_{|p|\to\infty}\sup_{\Omega\times [m,M]}\ \alpha,\,\beta,\,\gamma
\end{equation}

Then (15.28) follows from the Riccati inequality

\begin{equation}
\tag{15.34}
\chi'(z)+a\chi^{2}(z)+b\chi(z)+c+\varepsilon\le 0
\end{equation}

on \([m,M]\), with \(\chi=\omega\circ\psi^{-1}\) and \(\varepsilon>0\).
It has increasing auxiliary solutions when \(a\le0\), \(c\le0\), or
\(b\le-2\sqrt{|ac|}\); arbitrary \(a,b,c\) require sufficiently small
\(\operatorname{osc}_\Omega u\).

\begin{theorem}[Global gradient bound under the general structure conditions]\label{gt:thm-15-2}\dcref{ch15:15.2}\group{ch15-global-interior-gradient}
\source{gilbarg-trudinger:page-0379}
Let $u\in C^2(\Omega)\cap C^1(\bar{\Omega})$ satisfy (15.1) in the bounded domain $\Omega$. Suppose $Q$ is elliptic in $\Omega$ and there are scalar multipliers $r,s$ satisfying (15.32), together with either $a\leq0$ or $c\leq0$, where $a,c$ are defined by (15.27) and (15.33). Then

\begin{equation}
\tag{15.35}
\sup_{\Omega} |Du| \leq C
\end{equation}

where $C$ depends on the quantities in (15.32), $\operatorname{osc}_{\Omega}u$, and $\sup_{\partial\Omega}|Du|$.
\end{theorem}
\begin{proof}
\sketch Apply the maximum principle to the transformed gradient
$\bar v=|D(\psi^{-1}\circ u)|^2$. The structural bounds reduce the zero-order term to
$a\chi^2+b\chi+c+\chi'$. Standard increasing Riccati choices make this term
strictly negative when $a\leq0$ or $c\leq0$, and (15.29) converts the resulting bound to (15.35).
\end{proof}

(i) For uniformly elliptic \(Q\), take \(a_*^{ij}=a^{ij}\), \(c_i=0\),
\(r=s=0\). In addition to \(a^{ij}=O(\lambda)\) and
\(b=O(\lambda|p|^2)\), assume

\begin{equation}
\tag{15.36}
\delta a^{ij}=o(\lambda), \qquad \delta b=o(\lambda|p|^2) \qquad \text{as } |p|\to\infty
\end{equation}

Then \(c\le0\), so Theorem 15.2 gives the global gradient bound.

(ii) The prescribed mean curvature equation is

\begin{equation}
\tag{15.37}
\Delta u-\frac{D_i u\, D_j u}{(1+|Du|^2)}\,D_{ij}u-nH(x)\sqrt{1+|Du|^2}=0
\end{equation}


For \(H\in C^1(\overline\Omega)\), take

\[
a_{*}^{ij}=\delta^{ij}, \qquad c_i=-\frac{p_i}{1+|p|^{2}}, \qquad r=-1 \quad \text{and} \quad s=0
\]

which gives

\[
\alpha=-1+\frac{2}{1+|p|^{2}}, \qquad \beta=\frac{n H(x)\sqrt{1+|p|^{2}}}{|p|^{2}},
\]

\[
\gamma=-\frac{n p_i D_i H \, (1+|p|^{2})^{3/2}}{|p|^{4}}
\]

and hence

\[
a=-1, \qquad b=0, \qquad c=n \sup_{\Omega} |DH|.
\]

Thus \(a=-1\), and (15.29) gives

\[
\tag{15.38}
\sup_{\Omega} |Du| \le c_1 + c_2 n \sup_{\Omega} |H| + c_3 \sup_{\partial\Omega} |Du| \exp\!\bigl(c_4 n \sup_{\Omega} |DH| \, \osc_{\Omega} u\bigr)
\]

for constants \(c_1,c_2,c_3,c_4\).

(iii) If \(a_*^{ij}=\delta^{ij}\), as for (15.3) or in dimension two, then

\[
\alpha=\frac{1}{\mathscr{E}}(\delta+r)\mathscr{E}+\frac{(r+1)^{2}|p|^{2}}{4}\frac{1}{\mathscr{E}},
\]

\[
\tag{15.39}
\beta=\frac{1}{\mathscr{E}}\bigl[(\delta+s)\mathscr{E}+(\delta+r)b\bigr]+\frac{(r+1)s|p|^{2}}{2}\frac{1}{\mathscr{E}},
\]

\[
\gamma=\frac{1}{\mathscr{E}}(\delta+s)b+\frac{s^{2}}{4}\frac{|p|^{2}}{\mathscr{E}}.
\]

For \(r=-1\), \(s=0\),

\[
\alpha=\frac{1}{\mathscr{E}}(\delta-1)\mathscr{E},
\]

\[
\beta=\frac{1}{\mathscr{E}}\bigl[\delta\mathscr{E}+(\delta-1)b\bigr],
\]

\[
\gamma=\frac{1}{\mathscr{E}}\delta b.
\]


The same formulas hold whenever

\[
a_{*}^{ij}(p)=a_{*}^{ij}(p/|p|)
\]

since then \(\delta a_*^{ij}=0\).

\bigskip

\section*{15.3. Interior Gradient Bounds}

Let \(B=B_R(y)\Subset\Omega\) and choose the cutoff

\[
\tag{15.40}
\eta(x)=\left(1-\frac{|x-y|^{2}}{R^{2}}\right)^{\beta},\qquad \beta \ge 1.
\]

Set

\[
w=\eta \bar{v}
\]

Then \(w(y)=\bar v(y)\), \(w|_{\partial B}=0\). Multiplying (15.24) by
\(\eta\) and substituting \(\bar v=w/\eta\) gives

\begin{equation}\tag{15.41}
\begin{aligned}
&-\eta a_{*}^{ij}D_{ik}\bar{u}D_{jk}w+\frac{1}{2}a^{ij}D_{ij}w
+\frac{1}{2}\left(B_{i}-\frac{2}{\eta}a^{ij}D_{j}\eta\right)D_{i}w\\
&\quad+\psi' w[\omega(\delta+r+1)+(\delta+s)]a_{*}^{ij}D_{ij}\bar{u}\\
&\quad+\left\{\frac{\omega'}{\psi'}\mathscr{E}+\omega^{2}(\delta+r)\mathscr{E}
+\omega[(\delta+s)\mathscr{E}+(\delta+r)b]+(\delta+s)b\right.\\
&\qquad\left.+\frac{1}{\eta^{2}}a^{ij}D_{i}\eta D_{j}\eta
-\frac{1}{2\eta}a^{ij}D_{ij}\eta-\frac{1}{2\eta}B_{i}D_{i}\eta\right\}w=0.
\end{aligned}
\end{equation}

where \(B_i\) is given by (15.26). Introducing scalar multipliers \(t_i\)
and using (15.17) gives

\begin{equation}
\begin{aligned}
&-\eta a_*^{ij}D_i\bar u\,D_j\bar u+\frac12 a^{ij}D_{ij}w+\frac12 \widetilde B_iD_iw\\
&\quad +\psi'w[\omega(\delta+r+1)+(\delta+s)+(\tilde\delta+t)]a_*^{ij}D_{ij}\bar u\\
&\quad +\Bigl\{\frac{\omega'}{\psi'}\mathcal E+\omega^2(\delta+r)\mathcal E+\omega[(\delta+\tilde\delta+s+t)\mathcal E+(\delta+r)b\\
&\qquad\qquad -\frac{v}{2\eta}D_i\eta(\delta+r+1)c_i]+(\delta+\tilde\delta+s+t)b\\
&\qquad\qquad -\frac{v}{2\eta}D_i\eta(\delta+\tilde\delta+s+t+1)c_i\\
&\qquad\qquad +\frac{1}{\eta^2}a^{ij}D_i\eta D_j\eta-\frac{1}{2\eta}a^{ij}D_{ij}\eta\Bigr\}w=0
\end{aligned}
\tag{15.42}
\end{equation}

where

\begin{equation}
\tilde\delta=-\frac{1}{2\eta}D_i\eta\,c_i,
\qquad
t=-\frac{t_i}{2\eta}D_i\eta
\tag{15.43}
\end{equation}

and

\[
\widetilde B_i=B_i-\frac{2}{\eta}a^{ij}D_j\eta-\frac{v}{2\eta}D_j\eta(\delta_j+t_j)c_i.
\]

The cutoff inequality is

\begin{equation}
a^{ij}D_{ij}w+\widetilde B_iD_iw+2\widetilde G\mathcal Ew\ge 0
\tag{15.44}
\end{equation}

with

\begin{equation}
\widetilde G=\frac{\omega'}{\psi'}+\tilde\alpha\omega^2+\beta\omega+\tilde\gamma,
\tag{15.45}
\end{equation}


and

\[
\tilde{\alpha}=\alpha.
\]

\[
\tilde{\beta}=\beta+\frac{1}{\mathscr{E}}\left\{-\frac{|p|^{2}}{2\eta}D_{i}\eta(\delta+r+1)c_{i}+(\delta+t)\mathscr{E}
\qquad\qquad\qquad\qquad+\frac{|p|^{2}}{2\lambda_{*}}\left[(\delta+r+1)a_{*}^{ij}\right]\left[(\tilde{\delta}+t)a_{*}^{ij}\right]\right\},
\]

\begin{equation}
\tag{15.46}
\tilde{\gamma}=\gamma+\frac{1}{\mathscr{E}}\left\{\frac{|p|^{2}}{4\lambda_{*}}\sum \left[(\delta+t)a_{*}^{ij}\right]^{2}+(\delta+t)b\right.
\qquad\qquad-\frac{v}{2\eta}D_{i}\eta(\delta+\tilde{\delta}+s+t+1)c_{i}+\frac{1}{\eta^{2}}a^{ij}D_{i}\eta D_{j}\eta
\qquad\qquad-\frac{1}{2\eta}a^{ij}D_{ij}\eta+\frac{|p|^{2}}{2\lambda_{*}}\left[(\delta+s)a_{*}^{ij}\right]\left[(\tilde{\delta}+t)a_{*}^{ij}\right]\Biggr\}.
\end{equation}

Assume, in addition to (15.32),

\[
|p|^{\theta}(\partial_{k}+t_{k})a_{*}^{ij}=O\!\left(\sqrt{\lambda_{*}\mathscr{E}}/|p|\right)
\qquad \text{as } |p|\to\infty,
\]

\begin{equation}
\tag{15.47}
|p|^{2\theta}\lambda_{*},\,|p|^{\theta}(\partial_{i}+t_{i})\mathscr{E},\,|p|^{\theta}(\partial_{i}+t_{i})b=O(\mathscr{E})
\qquad \text{as } |p|\to\infty,
\end{equation}

\[
|p|^{\theta}(\tilde{\delta}+r+1)c_{i},\,|p|^{\theta}(\delta+s+1)c_{i},\,|p|^{2\theta}(\partial_{i}+t_{i})c_{i}=O(\mathscr{E}/|p|^{2})
\]

\[
\qquad\qquad\text{as } |p|\to\infty,
\]

for \(i,j,k=1,\ldots,n\), some \(\theta>0\), uniformly for \((x,z)\in(m,M)\).

With (15.40), \(\beta=2/\theta\), and sufficiently large \(|Du|\),

\[
|\beta-\tilde{\beta}|\leq C\frac{|D\eta|}{\eta|Du|^{\theta}}\leq \frac{C\beta}{Rw^{1/\beta}},
\]

\[
|\gamma-\tilde{\gamma}|\leq C\left\{\frac{|D\eta|^{2}}{\eta^{2}|Du|^{2\theta}}+\frac{|D\eta|}{\eta|Du|^{\theta}}\right\}\leq
\left\{\frac{\beta^{2}}{R^{2}w^{2/\beta}}+\frac{\beta}{Rw^{1/\beta}}\right\},
\]

where \(C\) depends on \(n\), (15.32), and (15.47). Hence

\begin{equation}
\tag{15.48}
|Du(y)|\leq C\left(1+R^{-1/\theta}\right),
\end{equation}


when \(a\le0\), \(c\le0\), \(b\le-2\sqrt{ac}\), or the oscillation on
\(B_R(y)\) is sufficiently small; the weak form extends this to \(C^2\) solutions.

\begin{theorem}[Local gradient estimate under the general structure conditions]\label{gilbarg-trudinger.thm.15.3.local-gradient-structure}\dcref{ch15:15.3}\group{ch15-global-interior-gradient}
Let $u\in C^2(\Omega)$ satisfy (15.1), and suppose $Q$ is elliptic and admits multipliers
$r,s,t_i$ satisfying (15.32) and (15.47). If $a\leq0$, $c\leq0$, or
$b\leq-2\sqrt{ac}$, with $a,b,c$ as in (15.27), (15.33), then

\begin{equation}
\tag{15.49}
|Du(y)| \leq C(1 + R^{-1/\theta})
\end{equation}

for $y\in\Omega$ and $B=B_R(y)$. The constant depends on (15.32), (15.47),
$\osc_{B\cap\Omega}u$, and $\sup_{B\cap\Omega}|Du|$. For arbitrary $a,b,c$, the same
estimate holds when $R$ is sufficiently small in terms of $a,b,c$ and the modulus of
continuity of $u$ at $y$.
\end{theorem}
\begin{proof}
\sketch Apply the transformed-gradient calculation to $\eta|Du|^2$ with the cutoff
(15.40). Conditions (15.32), (15.47) control the cutoff error terms and reduce the
maximum-point inequality to the same quadratic as in Theorem 15.2, giving (15.48)
and hence (15.49). Small oscillation handles arbitrary $a,b,c$.
\end{proof}

Under the natural uniformly elliptic conditions, the oscillation estimate is:

\begin{lemma}[Oscillation estimate under quadratic gradient growth]\label{gilbarg-trudinger.lem.15.4.oscillation-quadratic-growth}\dcref{ch15:15.4}\group{ch15-global-interior-gradient}
Let \(u\in W^{2,n}(\Omega)\) satisfy

\begin{equation}
\tag{15.50}
|Lu| \leq \lambda(\mu_0|Du|^2 + f)
\end{equation}

in $\Omega$, where $Lu=a^{ij}D_{ij}u$ satisfies (9.47), $\mu_0\in\mathbb R$, and
$f\in L^n(\Omega)$. Then for $B_0=B_{R_0}(y)\subset\Omega$ and $R\leq R_0$,

\begin{equation}
\tag{15.51}
\osc_{B_R(y)} u \leq C\left(\frac{R}{R_0}\right)^\alpha \left(\osc_{B_0}u + R\|f\|_{n;\Omega}\right)
\end{equation}

where $C,\alpha$ depend only on $n$, $\Lambda/\lambda$, and $\mu_0M$, with
$M=\|u\|_{0;\Omega}$.
\end{lemma}

\begin{proof}
\sketch First suppose $u\geq0$ satisfies

\begin{equation}
\tag{15.52}
Lu \leq \lambda(\mu_0|Du|^2 + f)
\end{equation}

for \(\mu_0>0\), \(f\ge0\), and set

\[
w = \frac{1}{\mu_0}\left(1 - e^{-\mu_0 u}\right)
\]



Then

\[
Lw \leq f
\]

in \(\Omega\). The weak Harnack inequality applies because

\[
w \leq u \leq e^{\mu_0 M}w,
\]

the corresponding constants depend on \(\mu_0M\). Oscillation iteration gives
(15.51), after applying the argument to shifted positive and negative parts.
\end{proof}

For Theorem 15.5 assume

\begin{equation}\tag{15.53}
\begin{aligned}
\Lambda,(\delta+r+1)a^{ij},(\delta+s)a^{ij},
 |p|^\theta(\partial_k+t_k)a^{ij}&=O(\lambda),\\
b,\ |p|^\theta(\partial_i+t_i)b&=O(\lambda|p|^2),\\
(\delta+r)b,\ (\delta+s)b&\leq O(\lambda|p|^2).
\end{aligned}
\end{equation}

as \(|p|\to\infty\), for some \(\theta>0\).

\begin{theorem}[Interior and global natural-growth gradient bounds]\label{gilbarg-trudinger.thm.15.5.natural-growth-gradient-bounds}\dcref{ch15:15.5}\group{ch15-global-interior-gradient}\uses{gilbarg-trudinger.thm.15.3.local-gradient-structure}\uses{gilbarg-trudinger.lem.15.4.oscillation-quadratic-growth}
Let $u\in C^2(\Omega)$ satisfy (15.1), and suppose $Q$ is elliptic and satisfies
(15.53) with multipliers $r,s,t_i$. For $y\in\Omega$,

\begin{equation}\tag{15.54}
|Du(y)| \leq C\left(1 + d_y^{-1/\theta}\right)
\end{equation}

where $C$ depends on $n$, the quantities in (15.53), $\|u\|_{0;\Omega}$, and
$d_y=\dist(y,\partial\Omega)$. If $u\in C^2(\Omega)\cap C^1(\bar\Omega)$, then

\begin{equation}\tag{15.55}
|Du(y)| \leq C
\end{equation}

where $C$ also depends on $\sup_{\partial\Omega}|Du|$.
\end{theorem}
\begin{proof}
\sketch Lemma 15.4 supplies the modulus of continuity required in Theorem 15.3.
Take $R$ comparable to $d_y$ for the interior estimate. Up to the boundary, combine
the same argument with the boundary gradient bound.
\end{proof}

\bigskip

\section*{15.4. Equations in Divergence Form}

Let \(u\in C^2(\Omega)\) satisfy

\begin{equation}\tag{15.56}
Qu = \divop A(x,u,Du) + B(x,u,Du) = 0
\end{equation}

in \(\Omega\). Its derivatives satisfy

\begin{equation}\tag{15.57}
\int_\Omega \left(\bar a^{ij}D_jD_k u + f_k^i\right)D_i\zeta\,dx = 0 \qquad \forall \zeta \in C_0^1(\Omega),
\end{equation}




where

\[
\tag{15.58}
\bar{a}^{ij}(x)=D_{p_j}A^i(x,u(x),Du(x)),
\qquad
f^i_k(x)=\delta_k A^i(x,u(x),Du(x))+\delta^{ik}B(x,u(x),Du(x)),\qquad \delta_k=p_k D_z+D_{x_k}.
\]

Assume

\[
\bar{a}^{ij}(x,z,p)\xi_i\xi_j=D_{p_j}A^i(x,z,p)\xi_i\xi_j\ge \nu(|z|)(1+|p|)^\tau |\xi|^2
\tag{15.59}
\]

for all \(\xi\), where \(\tau\in\mathbb R\) and \(\nu>0\) is nonincreasing, and

\[
|p|\,D_z A,\ D_x A,\ B=o(|p|)(1+|p|)^\tau \qquad \text{as } |p|\to\infty.
\tag{15.60}
\]

For \(M=\sup_\Omega|u|\), \(M_1=\sup_\Omega|Du|\), apply Theorem 8.16
to (15.57) on
\(\widetilde\Omega=\{x:|Du(x)|>M_1/(2\sqrt n)\}\):

\[
\sup_{\tilde{\Omega}} |D_k u|\le \sup_{\partial\tilde{\Omega}} |D_k u|+C(2\sqrt{n})^\tau \|(1+|Du|)^{-\tau}f^i_k\|_{q}
\tag{15.61}
\]

for \(q>n\), with \(C=C(n,\nu(M),q,|\Omega|)\). Taking \(q=\infty\) gives

\[
M_1=\sup_{\Omega}|Du|\le C\bigl(\sup_{\partial\Omega}|Du|+\sigma(M_1)\bigr)
\tag{15.62}
\]

where \(\sigma(t)=o(t)\).

\medskip

\begin{theorem}[Global gradient bound in divergence form]\label{gilbarg-trudinger.thm.15.6.divergence-global-gradient}\dcref{ch15:15.6}\group{ch15-global-interior-gradient}
Let $u\in C^2(\Omega)$ satisfy (15.56) in the bounded domain $\Omega$, and assume
(15.59), (15.60). Then

\[
\sup_{\Omega}|Du|\le C\bigl(1+\sup_{\partial\Omega}|Du|\bigr)
\tag{15.63}
\]

where $C$ depends on $n$, $\tau$, $\nu(M)$, and the quantities in (15.60).
\end{theorem}
\begin{proof}
\sketch Apply Theorem 8.16 to the differentiated equations (15.57) on the region
where $|Du|$ is larger than half its maximum. Estimate (15.61) and the little-$o$
condition (15.60) give (15.62), which absorbs the large-gradient term and yields (15.63).
\end{proof}

For the interior estimate assume uniform ellipticity in the form

\begin{equation}
|D_p A(x,z,p)| \le \mu(|z|)(1+|p|)^\tau
\tag{15.64}
\end{equation}

where \(\mu>0\) is nondecreasing. If \(\tau>-1\), then (15.59), (15.64) imply

\[
p\cdot A(x,z,p)-p\cdot A(x,z,0)\ge \frac{\nu}{\tau+1}|p|^{\tau+2}.
\]

\begin{equation}
|A^i(x,z,p)-A^i(x,z,0)|\le \frac{\mu}{\tau+1}(1+|p|)^{\tau+1},\qquad i=1,\dots,n.
\tag{15.65}
\end{equation}

Assume also

\begin{equation}
g(x,z,p)=(1+|p|)\,|D_z A|+|D_x A|+|B|\le \mu(|z|)(1+|p|)^{\tau+2}
\tag{15.66}
\end{equation}

on \(\Omega\times\mathbb R\times\mathbb R^n\).

\medskip

\noindent\textit{(i) \(L^p\) reduction.}
Test (15.57) with \(\zeta D_ku\), sum in \(k\), set \(v=|Du|^2\), and
apply Young's inequality:

\begin{equation}
\int_{\Omega}\left(\tilde{a}^{ij}D_j v+2D_k u\,f_k^i\right)D_i\zeta\,dx\le \frac12\int_{\Omega}\zeta\lambda^{-1}\sum (f_k^i)^2\,dx
\tag{15.67}
\end{equation}

for nonnegative \(\zeta\in C_0^1(\Omega)\). Set

\[
\bar{v}=\int_0^v (1+\sqrt{t})^\tau\,dt,
\]

Then (15.67) becomes

\begin{equation}
\int_{\Omega}\left(\tilde{a}^{ij}D_j \bar{v}+2D_k u\,f_k^i\right)D_i\zeta\le \frac12\int_{\Omega}\zeta\lambda^{-1}\sum (f_k^i)^2\,dx,
\tag{15.68}
\end{equation}

where

\[
\tilde{a}^{ij}=\bar{a}^{ij}(1+|Du|)^{-\tau}.
\]


Theorem 8.17 and (15.66) imply, for sufficiently large \(p\),
\begin{equation}
\sup_{B_R(y)} v \le C(n,\nu(M),\mu(M),\tau,\operatorname{diam}\Omega,R^{-n}\int_{B_{2R}(y)} v^p\,dx).
\tag{15.69}
\end{equation}

\noindent\textit{(ii) H\"older reduction.}
The weak equation is
\begin{equation}
\mathcal{Q}(u,\varphi)=\int_{\Omega}\bigl(A^i(x,u,Du)D_i\varphi-B(x,u,Du)\varphi\bigr)\,dx=0 \qquad \forall \varphi\in C_0^1(\Omega).
\tag{15.70}
\end{equation}

By (15.65),
\begin{equation}
|A(x,z,p)|\le \mu_1(|z|)(1+|p|)^{\tau+1}
\tag{15.71}
\end{equation}
\[
p\cdot A(x,z,p)\ge \nu_1(|z|)|p|^{\tau+2}-\mu_1(|z|)
\]
for $(x,z,p)\in \Omega\times \mathbb{R}\times \mathbb{R}^n$, where $\mu_1$ and $\nu_1$ are respectively positive non-decreasing and positive non-increasing functions depending on $\tau,\mu,\nu$ and $\sup_{\Omega}|A(x,u,0)|$.

Test (15.70) with
\[
\varphi=\eta^2[u-u(y)]
\]
where \(\eta\in C_0^1(B_{2R}(y))\). With
\(\omega(R)=\sup_{B_{2R}}|u(x)-u(y)|\), choose \(R\) so that
\(\omega(R)\le\nu_1 2^{-\tau}/(8(\mu+\mu_1))\).
Then
\[
\int_{\Omega} \eta^2 |Du|^{\tau+2}\,dx \leq C \int_{\Omega} \bigl(\eta^2+\omega(R)(1+|Du|)^{\tau}|D\eta|^2\bigr)\,dx
\tag{15.72}
\]
where \(C=C(\mu,\mu_1,\nu_1,\tau)\). Testing with
\[
\eta v^{(\beta+1)/2}, \qquad \beta>0,
\]
gives
\[
\int_{\Omega} \eta^2 |Du|^{2\beta+\tau+4}\,dx \leq C \int_{\Omega} \Bigl\{ |Du|^{2(\beta+1)}\bigl(\eta^2+\omega(R)(1+|Du|)^{\tau}|D\eta|^2\bigr)
\qquad\qquad +(\beta+1)^2\omega(R)(1+|Du|)^{\tau}\eta^2 v^{\beta-1}|Dv|^2 \Bigr\}\,dx.
\tag{15.73}
\]
Testing (15.67) with \(\zeta=\eta^2v^\beta\), applying Young's inequality,
and taking \(\omega(R)\) sufficiently small gives a cutoff estimate. With
\(\eta=1\) on \(B_R(y)\) and
\(|D\eta|\le2/R\), gives for every \(p\ge1\)
\begin{equation}\tag{15.74}
\int_{B_R(y)}(1+|Du|)^p\,dx\leq C
\end{equation}
where \(C=C(\mu,\nu,\mu_1,\nu_1,p,\tau,R^{-1})\). Combining (15.69) and
(15.74), for \(B_0=B_{R_0}(y)\Subset\Omega\) and \(0<\alpha<1\),
\begin{equation}\tag{15.75}
|Du(y)|\leq C
\end{equation}
where \(C=C(n,\tau,\nu,\mu,\nu_1,\mu_1,\alpha,[u]_{\alpha,y})\) and
\[
[u]_{\alpha,y}=\sup_{B_0}\frac{|u(x)-u(y)|}{|x-y|^\alpha}.
\]

\noindent\textit{(iii) H\"older estimate for weak solutions.}
Rewrite (15.71), (15.66) as
\begin{equation}\tag{15.76}
\begin{aligned}
|A(x,z,p)|&\leq a_0|p|^{m-1}+\chi^{m-1} \\
p\cdot A(x,z,p)&\geq \nu_0|p|^m-\chi^m \\
|B(x,z,p)|&\leq b_0|p|^m+\chi^m
\end{aligned}
\end{equation}


where \(m=\tau+2>1\) and \(a_0,b_0,\nu_0,\chi>0\) may depend on
\(M=\sup_\Omega|u|\).

\begin{theorem}[H\"older estimate for weak divergence-form solutions]\label{gilbarg-trudinger.thm.15.7.weak-solution-holder}\dcref{ch15:15.7}\group{ch15-global-interior-gradient}
Let $u\in C^1(\Omega)$ be a weak solution of (15.56) and assume (15.76).
For $B_0=B_{R_0}(y)\subset\Omega$ and $R\leq R_0$,

\begin{equation}
\tag{15.77}
\osc_{B_R(y)} u \leq C(1+R_0^{-\alpha}M_0)R^\alpha
\end{equation}

where $C=C(n,a_0,b_0,\nu_0,\chi,R_0,m,M_0)$ and
$\alpha=\alpha(n,a_0,b_0,\nu_0,m,M_0)>0$, with $M_0=\sup_{B_0}|u|$.
\end{theorem}

\begin{proof}
\sketch For
\[
k=\chi(R+R^{m/(m-1)})
\]
the weak Harnack argument gives
\begin{equation}
\tag{15.78}
R^{-n/p}\|u\|_{L^p(B_{2R}(y))}\leq C\bigl(\inf_{B_R(y)} u+k\bigr)
\end{equation}
where \(C=C(n,a_0,b_0,\nu_0,m,M_0,p)\) and \(1\le p<n/(n-m)^+\). Use
the test function
\[
\varphi=\eta^m \bar{u}^{\beta} \exp(-b_0\bar{u}), \qquad \beta<0,
\]
in \(\mathcal Q(u,\varphi)\ge0\), followed by Sobolev (7.26) and oscillation
iteration (Theorem 8.22), to obtain (15.77).
\end{proof}

\begin{theorem}[Interior gradient bound in divergence form]\label{gilbarg-trudinger.thm.15.8.divergence-interior-gradient}\dcref{ch15:15.8}\group{ch15-global-interior-gradient}\uses{gilbarg-trudinger.thm.15.7.weak-solution-holder}
Let $u\in C^2(\Omega)$ satisfy (15.57), and assume (15.60), (15.64), (15.66)
with $\tau>-1$. Then

\begin{equation}
\tag{15.79}
|Du(y)|\leq C
\end{equation}


for $y\in\Omega$, where $d=\dist(y,\partial\Omega)$,
$M_0=\sup_{B_d(y)}|u|$, and $C$ depends on
$n,\tau,\nu(M_0),\mu(M_0),\sup_\Omega|A(x,u,0)|,M_0,M_0/d$.
\end{theorem}
\begin{proof}
\sketch The energy estimate (15.68) reduces the gradient bound to an $L^p$ estimate
(15.69). The H\"older control from Theorem 15.7 yields the integral estimate (15.74);
substitution in (15.69) gives (15.75), hence (15.79).
\end{proof}

\begin{theorem}[Global divergence-form gradient bound]\label{gilbarg-trudinger.thm.15.9.divergence-global-bound}\dcref{ch15:15.9}\group{ch15-global-interior-gradient}\uses{gilbarg-trudinger.thm.15.8.divergence-interior-gradient}
Let $u\in C^2(\Omega)\cap C^0(\overline\Omega)$ satisfy (15.57) in bounded
$\Omega$, and assume (15.60), (15.64), (15.66) with $\tau>-1$. Suppose $\Omega$
satisfies a uniform exterior sphere condition and $u=\varphi$ on $\partial\Omega$ for
$\varphi\in C^2(\overline\Omega)$. Then

\[
\tag{15.80}
\sup_{\Omega} |Du| \leq C
\]

where $C$ depends on $n,\tau,\nu(M),\mu(M),\sup_\Omega|A(x,u,0)|$,
$\partial\Omega$, $|\varphi|_2$, and $M=\sup_\Omega|u|$.
\end{theorem}
\begin{proof}
\sketch Combine Theorem 15.8 on interior balls with the boundary Lipschitz estimate
of Theorem 14.1, using the exterior-sphere geometry to cover the boundary layer.
\end{proof}

\medskip

\section*{15.5. Existence Theorems}

\subsection*{(i) Uniformly elliptic equations in the general form (15.1)}
Assume

\begin{equation}\tag{15.81}
\begin{aligned}
a^{ij},\delta a^{ij}&=O(\lambda),&
\delta a^{ij}&=o(\lambda),\\
b&=O(\lambda|p|^2),&
\delta b&\le O(\lambda|p|^2),\\
\delta b&\le o(\lambda|p|^2).
\end{aligned}
\end{equation}

as \(|p|\to\infty\), uniformly for \(x\in\Omega\) and bounded \(z\).

\begin{theorem}[Uniformly elliptic Dirichlet existence]\label{gilbarg-trudinger.thm.15.10.uniformly-elliptic-existence}\dcref{ch15:15.10}\group{ch15-global-interior-gradient}
Let $\Omega$ be bounded and $Q$ elliptic in $\overline\Omega$, with
$a^{ij},b\in C^1(\overline\Omega\times\mathbb R\times\mathbb R^n)$ satisfying
(15.81) or (15.53) and (10.10) (or (10.36)). If
$\partial\Omega\in C^{2,\gamma}$ and $\varphi\in C^{2,\gamma}(\overline\Omega)$,
$0<\gamma<1$, then the Dirichlet problem $Qu=0$, $u|_{\partial\Omega}=\varphi$
has a solution $u\in C^{2,\gamma}(\overline\Omega)$.
\end{theorem}
\begin{proof}
\sketch The global gradient estimate and the height estimate from the Chapter 10
maximum principle give the a priori bounds required by the continuity method.
Schauder estimates provide compactness and solvability.
\end{proof}

If the maximum-principle hypothesis is omitted, solvability instead requires
uniform boundedness of the associated family (13.42). Uniqueness follows when
\(a^{ij}\) is independent of \(z\) and \(D_zb\le0\).

\medskip

\subsection*{(ii) Uniformly elliptic divergence-form equations (15.56)}
Assume

\begin{equation}\tag{15.82}
\begin{aligned}
|p|^\tau &\leq O(\lambda)\\
D_p A &= O(|p|^\tau),\\
|p|\,D_z A,\ D_x A,\ B &= O(|p|^{\tau+2})
\end{aligned}
\end{equation}

as \(|p|\to\infty\), uniformly for \(x\in\Omega\) and bounded \(z\), with \(\tau>-1\).

\medskip

\begin{theorem}[Uniformly elliptic divergence-form Dirichlet existence]\label{gilbarg-trudinger.thm.15.11.divergence-existence}\dcref{ch15:15.11}\group{ch15-global-interior-gradient}
Let \(\Omega\subset\mathbb R^n\) be bounded and let \(Q\) be elliptic in
\(\overline\Omega\), with \(A^i\in C^{1,\gamma}(\overline\Omega\times\mathbb R\times\mathbb R^n)\),
\(B\in C^\gamma(\overline\Omega\times\mathbb R\times\mathbb R^n)\), \(0<\gamma<1\).
Assume (15.82), the hypotheses of Theorem 10.9 with \(\alpha=\tau+2\), and
\(\partial\Omega,\varphi\in C^{2,\gamma}\). Then the Dirichlet problem
\(Qu=0\) in \(\Omega\), \(u=\varphi\) on \(\partial\Omega\), has a solution
\(u\in C^{2,\gamma}(\overline\Omega)\).
\end{theorem}
\begin{proof}
\sketch Use the homotopy (15.83). The height estimate from Theorem 10.9 and
the boundary and interior gradient estimates (14.1), (15.9) give uniform
\(C^1\) bounds along the family; Schauder compactness and the continuity method
then yield a solution at \(\sigma=1\).
\end{proof}

\medskip

Use the homotopy

\begin{equation}\tag{15.83}
Q_\sigma u=\diver\{\sigma A+(1-\sigma)(1+|Du|^2)^{\tau/2}Du\}+\sigma B=0,
\end{equation}

\[
u=\sigma\varphi \text{ on } \partial\Omega,\qquad 0\leq \sigma \leq 1.
\]

If the maximum-principle hypotheses are unavailable, require uniform boundedness
of this homotopy. Equations of the form

\begin{equation}\tag{15.84}
Qu=a^{ij}(x,u)D_{ij}u+b(x,u,Du)
\end{equation}

with \(a^{ij}\in C^1(\Omega\times \R)\) can be written in the divergence form (15.57) with

\[
A^i(x,z,p)=a^{ij}(x,z)p_j,
\]
\[
B(x,z,p)=b(x,z,p)-D_z a^{ij}(x,z)p_ip_j-D_{x_i}a^{ij}(x,z)p_j.
\]


\medskip

\begin{theorem}[Nondivergence Dirichlet existence via divergence form]\label{gilbarg-trudinger.thm.15.12.nondivergence-existence}\dcref{ch15:15.12}\group{ch15-global-interior-gradient}
Let \(\Omega\) be bounded and let \(Q\) be the elliptic operator (15.84) on
\(\overline\Omega\), with \(a^{ij}\in C^1(\overline\Omega\times\mathbb R)\) and
\(b\in C^\gamma(\overline\Omega\times\mathbb R\times\mathbb R^n)\), \(0<\gamma<1\).
Assume \(b=O(|p|^2)\), (10.10) or (10.36), and
\(\partial\Omega,\varphi\in C^{2,\gamma}\). Then \(Qu=0\) in \(\Omega\),
\(u=\varphi\) on \(\partial\Omega\), has a solution in
\(C^{2,\gamma}(\overline\Omega)\).
\end{theorem}
\begin{proof}
\sketch Rewrite (15.84) in divergence form (15.57) using the displayed
\(A^i\) and \(B\). The estimates of Theorems 10.3, 14.1, and 15.9 provide
\(C^1\) a priori bounds, after which the continuity method gives the stated
\(C^{2,\gamma}\) solution.
\end{proof}

\medskip

\subsection*{(iii) Non-uniformly elliptic equations in general domains}
Assume the coefficients in (15.1) decompose as in (15.23) and satisfy
the structure conditions

\begin{equation}\tag{15.85}
\begin{aligned}
|p|\, a^{ij},\, b &= O(\mathcal{E}), \\
\delta a^{ij}_* &= O\!\left(\sqrt{\lambda_*\,\mathcal{E}/|p|}\right), \\
\delta a_*^{ij} &= o\!\left(\sqrt{\lambda_*\,\mathcal{E}/|p|}\right), \\
\delta b &\leq O(\mathcal{E}), \\
\delta b &\leq o(\mathcal{E}),
\end{aligned}
\end{equation}

as \(|p|\to\infty\), uniformly for \(x\in\Omega\) and bounded \(z\).

\medskip

\begin{theorem}[General-domain nonuniform elliptic existence]\label{gilbarg-trudinger.thm.15.13.general-domain-existence}\dcref{ch15:15.13}\group{ch15-global-interior-gradient}
Let \(\Omega\subset\mathbb R^n\) be bounded, with
\(a^{ij},b\in C^1(\overline\Omega\times\mathbb R\times\mathbb R^n)\) satisfying
(15.85) and (10.10) or (10.36). If
\(\partial\Omega,\varphi\in C^{2,\gamma}\), \(0<\gamma<1\), then the
Dirichlet problem \(Qu=0\), \(u|_{\partial\Omega}=\varphi\), has a solution
\(u\in C^{2,\gamma}(\overline\Omega)\).
\end{theorem}
\begin{proof}
\sketch Combine the global gradient estimate from Theorem 15.2 with the
boundary estimate of Theorem 14.1 and the height bound from Chapter 10.
The continuity method and Schauder estimates then give the solution.
\end{proof}

\medskip

When (15.2) holds, \(\delta a_*^{ij}=0\); the same conclusion holds when
\(a_*^{ij}(p)=a_*^{ij}(p/|p|)\).

\medskip

\subsection*{(iv) Non-uniformly elliptic equations in convex domains}
Assume the decomposition (15.2) and

\begin{equation}\tag{15.86}
\begin{aligned}
b &= o(|p|\Lambda), \\
\delta b &\leq O\!\left(b^2/|p|^2\,\mathcal{T}_*\right), \qquad (\mathcal{T}_* = \operatorname{trace}[a_*^{ij}]),
\end{aligned}
\end{equation}

as \(|p|\to\infty\), uniformly in \(x\in\Omega\) and bounded \(z\).

\medskip

\begin{theorem}[Uniformly convex-domain existence]\label{gilbarg-trudinger.thm.15.14.convex-domain-existence}\dcref{ch15:15.14}\group{ch15-global-interior-gradient}
Let \(\Omega\subset\mathbb R^n\) be uniformly convex and let \(Q\) be elliptic
with \(a^{ij},b\in C^1(\overline\Omega\times\mathbb R\times\mathbb R^n)\)
satisfying (15.86). If


\(\partial\Omega,\varphi\in C^{2,\gamma}\), \(0<\gamma<1\), then the
Dirichlet problem \(Qu=0\), \(u|_{\partial\Omega}=\varphi\), is solvable provided
the family of \(C^{2,\gamma}(\overline\Omega)\) solutions of (15.87), for some
fixed \(x_0\in\Omega\), is uniformly bounded:

\[
\tag{15.87}
Q_\sigma u = a^{ij}(x,u,Du)D_{ij}u + \sigma b(\sigma x + (1-\sigma)x_0,\sigma u,Du)=0,
\]
\[
u=\sigma\varphi \text{ on } \partial \Omega,\qquad 0\le \sigma \le 1,
\]

\end{theorem}
\begin{proof}
\sketch The gradient maximum principle from Theorem 15.1 and the boundary
estimate of Corollary 14.5 give uniform estimates along (15.87); the assumed
family bound supplies the remaining a priori control for the continuity method.
\end{proof}

\medskip

For the divergence-form equation (15.57), assume for some \(\tau\in\mathbb R\),
\begin{equation}\tag{15.88}
|p|^\tau\le O(\lambda),\qquad
|p|\,D_zA,\ D_xA,\ B=o(|p|^{\tau+1})
\end{equation}
as \(|p|\to\infty\), uniformly for \(x\in\Omega\) and bounded \(z\).

\medskip

\begin{theorem}[Convex-domain divergence-form existence]\label{gilbarg-trudinger.thm.15.15.convex-divergence-existence}\dcref{ch15:15.15}\group{ch15-global-interior-gradient}
Let \(\Omega\subset\mathbb R^n\) be uniformly convex and let \(Q\) be elliptic
on \(\overline\Omega\), with
\(A^i,B\in C^1(\overline\Omega\times\mathbb R\times\mathbb R^n)\) satisfying
(15.88). If \(\partial\Omega,\varphi\in C^{2,\gamma}\), \(0<\gamma<1\),
then \(Qu=0\) in \(\Omega\), \(u=\varphi\) on \(\partial\Omega\), has a
solution \(u\in C^{2,\gamma}(\overline\Omega)\), provided the
\(C^{2,\gamma}(\overline\Omega)\) solutions of (13.42) are uniformly bounded.
\end{theorem}
\begin{proof}
\sketch Theorem 15.6 supplies the global gradient bound for the divergence-form
homotopy, while the convex-boundary estimate controls the boundary gradient.
Uniform height bounds and the continuity method give the solution.
\end{proof}

\medskip

\subsection*{(v) Boundary curvature conditions}
Assume the decompositions (14.43) and (15.23):

\[
a^{ij}(x,z,p)=\Lambda a_\infty^{ij}(x,p/|p|)+a_0^{ij}(x,z,p)
\]
\[
\tag{15.89}
= a_*^{ij}(x,z,p)+\frac12\bigl[p_ic_j(x,z,p)+c_i(x,z,p)p_j\bigr],
\]
\[
b(x,z,p)=|p|\Lambda b_\infty(x,z,p/|p|)+b_0(x,z,p)
\]

\noindent where $a_\infty^{ij}\in C^1(\overline{\Omega}\times B(0))$, $a_*^{ij}, b_\infty\in C^1(\overline{\Omega}\times\mathbb{R}\times B(0))$, $i,j=1,\ldots,n$, the matrices $[a_\infty^{ij}]$, $[a_*^{ij}]$ are non-negative and symmetric, and $b_\infty$ is non-increasing with respect to $z$. Impose

\[
\delta\mathcal{E}\le \mathcal{E}+o(\mathcal{E}),
\]
\[
\delta\mathcal{E},\,\delta b,\;(\delta-1)b_0\le O(\mathcal{E}),
\]
\[
\tag{15.90}
\delta a_*^{ij}=O\!\left(\frac{\sqrt{\mathcal{E}}}{|p|}\right),
\]
\[
a_0^{ij}=o(\Lambda)
\]
\[
b_0=o(|p|)
\]



as \(|p|\to\infty\), uniformly for \(x\in\Omega\) and bounded \(z\).

\medskip

\begin{theorem}[Existence under strict boundary curvature inequalities]\label{gilbarg-trudinger.thm.15.16.strict-boundary-curvature}\dcref{ch15:15.16}\group{ch15-global-interior-gradient}
Let \(\Omega\subset\mathbb R^n\) be bounded and let \(Q\) be elliptic on
\(\overline\Omega\), with
\(a^{ij},b\in C^1(\overline\Omega\times\mathbb R\times\mathbb R^n)\) satisfying
(15.89), (15.90). Suppose \(\partial\Omega,\varphi\in C^{2,\gamma}\),
\(0<\gamma<1\), and

\[
\tag{15.91}
\mathcal{K}^{+}>b_\infty(y,\varphi(y),v), \qquad \mathcal{K}^{-}>-b_\infty(y,\varphi(y),-v)
\]

at every \(y\in\partial\Omega\), where \(v\) is the inward unit normal and
\(\mathcal K^\pm\) are the nonnegative quantities from (14.44). Then the
Dirichlet problem \(Qu=0\), \(u|_{\partial\Omega}=\varphi\), is solvable
provided the \(C^{2,\gamma}(\overline\Omega)\) solutions of (13.42) are
uniformly bounded.
\end{theorem}
\begin{proof}
\sketch The strict inequalities (15.91) give the boundary gradient barriers
of Theorem 14.9. Combine them with Theorem 15.2 and the assumed height bound
along (13.42), then apply the continuity method.
\end{proof}

\medskip

For non-strict boundary inequalities, strengthen (15.90) to

\[
\delta\mathcal{E}\le \mathcal{E}+o(\mathcal{E}),
\]
\[
\delta\mathcal{E},\,\delta b,\,\bar{\delta}b_0\le O(\mathcal{E}),
\]
\[
\tag{15.92}
\delta a_*^{ij}=O\!\left(\sqrt{\mathcal{E}}/|p|\right),
\]
\[
a_0^{ij}=O(\mathcal{E}/|p|),
\]
\[
b_0=O(\mathcal{E}),
\]

\noindent as \(|p|\to\infty\), uniformly for \(x\in\Omega\) and bounded \(z\).

\medskip

\begin{theorem}[Existence under non-strict boundary curvature inequalities]\label{gilbarg-trudinger.thm.15.17.nonstrict-boundary-curvature}\dcref{ch15:15.17}\group{ch15-global-interior-gradient}
Let \(\Omega\subset\mathbb R^n\) be bounded and let \(Q\) be elliptic on
\(\overline\Omega\), with
\(a^{ij},b\in C^1(\overline\Omega\times\mathbb R\times\mathbb R^n)\) satisfying
(15.89), (15.92). Suppose \(\partial\Omega,\varphi\in C^{2,\gamma}\),
\(0<\gamma<1\), and

\[
\tag{15.93}
\mathcal{K}^{+}\ge b_\infty(y,\varphi(y),v), \qquad \mathcal{K}^{-}\ge -b_\infty(y,\varphi(y),-v), \qquad \mathcal{K}^{\pm}\ge 0
\]

at every \(y\in\partial\Omega\). Then \(Qu=0\) in \(\Omega\),
\(u=\varphi\) on \(\partial\Omega\), is solvable provided the
\(C^{2,\gamma}(\overline\Omega)\) solutions of (13.42) are uniformly bounded.
\end{theorem}
\begin{proof}
\sketch The strengthened conditions (15.92) permit the boundary estimate of
Theorem 14.6 under the non-strict inequalities (15.93). Theorem 15.2 and the
uniform family bound complete the a priori estimates for the continuity method.
\end{proof}

\bigskip

\section*{15.6. Existence for Continuous Boundary Values}

\medskip

Approximate \(\varphi\in C^0(\partial\Omega)\) uniformly by
\(\varphi_m\in C^{2,\gamma}\). Interior estimates (15.3), (15.6), (13.1),
(13.3), and (6.2), together with the boundary modulus estimate (14.15), give
compactness of the corresponding solutions and convergence to the continuous
boundary data. Uniform height bounds are supplied by the relevant maximum
principle.

\medskip

\begin{theorem}[Continuous-data existence for nondivergence equations]\label{gilbarg-trudinger.thm.15.18.continuous-data}\dcref{ch15:15.18}\group{ch15-global-interior-gradient}
Let \(\Omega\subset\mathbb R^n\) be bounded and satisfy an exterior sphere
condition at every boundary point. Let \(Q\) be elliptic on \(\Omega\), with
\(a^{ij},b\in C^1(\Omega\times\mathbb R\times\mathbb R^n)\), satisfying the
hypotheses of Theorem 15.3 or 15.5, Theorem 14.1, and (10.10) or (10.36).
For every \(\varphi\in C^0(\partial\Omega)\), the problem
\(Qu=0\) in \(\Omega\), \(u=\varphi\) on \(\partial\Omega\), has a solution
\(u\in C^0(\overline\Omega)\cap C^2(\Omega)\).
\end{theorem}
\begin{proof}
\sketch Solve the smooth approximating boundary problems, use the interior
compactness and boundary modulus estimate (14.15), and pass to the limit.
\end{proof}

\medskip

\begin{theorem}[Continuous-data existence in divergence form]\label{gilbarg-trudinger.thm.15.19.continuous-divergence-data}\dcref{ch15:15.19}\group{ch15-global-interior-gradient}
Let \(\Omega\subset\mathbb R^n\) be bounded and satisfy an exterior sphere
condition at every boundary point. Let \(Q\) be a divergence-form operator
with \(A^i\in C^{1,\gamma}(\Omega\times\mathbb R\times\mathbb R^n)\),
\(B\in C^\gamma(\Omega\times\mathbb R\times\mathbb R^n)\), \(0<\gamma<1\).
Assume the hypotheses of Theorem 15.8 and Theorem 10.9 with
\(\alpha=\gamma+2\). For every \(\varphi\in C^0(\partial\Omega)\), the
problem \(Qu=0\) in \(\Omega\), \(u=\varphi\) on \(\partial\Omega\), has a
solution \(u\in C^0(\overline\Omega)\cap C^{2,\gamma}(\Omega)\).
\end{theorem}
\begin{proof}
\sketch Approximate the boundary data, apply the divergence-form interior
estimate (15.8) and Theorem 10.9, and use interior Schauder compactness plus
the boundary modulus estimate (14.15) to pass to the limit.
\end{proof}

\bigskip
